\documentclass[11pt,a4paper]{article}
\usepackage[T1]{fontenc}
\usepackage[utf8]{inputenc}
\usepackage{lmodern,microtype}
\usepackage[a4paper,top=24mm,bottom=24mm,left=25mm,right=25mm,headheight=14pt]{geometry}
\usepackage{amsmath,amssymb,amsthm,mathtools,booktabs,array,tabularx,xcolor,fancyhdr}
\usepackage[hyphens]{url}
\usepackage[colorlinks=true,linkcolor=blue!30!black,urlcolor=blue!45!black,citecolor=blue!45!black]{hyperref}
\definecolor{ink}{HTML}{223D50}
\definecolor{accent}{HTML}{176459}
\pagestyle{fancy}\fancyhf{}\fancyhead[L]{\small\sffamily\color{ink} Reflected Hankel determinants}\fancyhead[R]{\small\sffamily\color{ink} Research manuscript}\fancyfoot[C]{\small\thepage}
\renewcommand{\headrulewidth}{0.3pt}
\allowdisplaybreaks[1]\setlength{\emergencystretch}{2em}
\hypersetup{pdftitle={Reflection formulas for Catalan Hankel determinants and a prime-cube evaluation},pdfauthor={Keio research series}}
\begin{document}
\begin{center}
{\small\sffamily\color{accent} RESEARCH MANUSCRIPT}\par\vspace{7pt}
{\LARGE\bfseries\color{ink} Reflection formulas for Catalan Hankel determinants and a prime-cube evaluation}\par\vspace{7pt}
{\small Keio research series\quad$\cdot$\quad 20 November 2026}
\end{center}
\begin{abstract}
We prove the five identities of Cigler's 2018 Conjecture 7.2 for every odd convolution power of the Catalan sequence. In the last identity two unbounded-order determinants have a sum of constant absolute value. Reflecting their frequency indices makes all but one column of a Taylor-coefficient determinant agree; a sine subtraction cancels the dependence on the order in the last column. It also determines the full validity range of the companion even conjecture. For the ordinary binomial moments, we prove both parities of Wang and Xin's Conjecture 22 on exact degrees and evaluate the two adjacent sums posed by Cigler in 2018. We compute the signed leading coefficients in Cigler's 2023 Conjectures 8, 9 and 13 by positive moment determinants, determine their least common denominators, correct a Bernoulli sign in Conjecture 14, and classify all five shifted binomial identities of Conjecture 17. An independent arithmetic result gives a congruence modulo $p^3$ for a cubic binomial convolution with a $p$-adic parameter. At parameter $1/4$ it gives a counterexample to one clause of Sun's Conjecture 35(v) for every prime $p>3$ congruent to $3$ modulo $8$.
\end{abstract}

\section{Problems and precedents}
Let $H_r(N)$ be the order-$N$ Hankel determinant of the coefficients of $C(z)^r$, where $C(z)$ is the Catalan generating function. A principal result of this paper is
\[
 \boxed{\ H_{2k+1}((2k+1)n-1)+H_{2k+1}((2k+1)n+2)
      =(-1)^{kn+1}(k-1)(2k+1)\ }\qquad(k,n\ge1).
\]
For example, at $k=3,n=1$ the two determinants are $329$ and $-315$. The same cancellation persists for every $n$. This is the last line of Cigler's Conjecture~7.2, equation (7.6), as posed in 2018 \cite{C18}. Wang and Xin subsequently studied these determinants by continued fractions at bounded convolution orders \cite{WX}. Cigler and Tyson recorded related Catalan identities and evaluated several unshifted binomial determinants \cite{CT}. For the different input $\binom{2h+r}{h}$ we prove Wang and Xin's Conjecture 22 in the numbering of their preprint: the four odd-order subsequences and the two signed even-order subsequences have exactly the degrees proposed in \cite{WX}. The 2018 Catalan problem and the 2023 numbered conjectures \cite{C23} use different numbering.

\paragraph{What changes beyond the degree calculation.}
Equality of leading coefficients in the boxed identity would still leave every lower-degree term to be accounted for. We compare entire columns before taking any asymptotic limit. For the 2023 shift question we must also recover the orientation and scale of the root-evaluation block. A positive Gram determinant gives the surviving coefficient in the multiple-root block. We further show that the normalized leading coefficients have an optimal common denominator. Specifically, for even order $2k$ and odd order $2k+1$ the respective normalized families each have least common denominator $\prod_{i=1}^{k-1}(2i-1)!!$. An arithmetic application of a different argument is given at the end. There the complementary index ranges factor a triangular convolution, and a second-order calculation determines a residue that a prime-square estimate does not detect.

\paragraph{Relation to the corpus.}
The motivating snapshot of the Verifiable Mathematical Reasoning Corpus ends at manuscript 225 \cite{VMRC}. Its manuscript 079 uses finite orthogonal-polynomial modifications of Hankel moments, and manuscript 026 uses positive moment expansions to identify a leading term. We use those ideas together with the classical determinantal framework \cite{Krat}. Earlier work in this series obtained the polynomial degrees by folding evaluation columns. Here the exact scalar and signs are retained, and the reflection of an entire pair of matrices yields an exact sum. The proofs are included below. The Christoffel identity is classical, and Cigler discusses the trigonometric expansions explicitly \cite{C23,MO}.

\section{A determinant problem with two reflected groups}
Let $C(z)=1+zC(z)^2$ with $C(0)=1$. Throughout the first part put
\[
 c_{r,i}=[z^i]C(z)^r=\frac{r}{2i+r}\binom{2i+r}{i}\quad(i\ge0),\qquad
 D_{r,s}(N)=\det(c_{r,i+j+s})_{0\le i,j<N}.
\]
Set $c_{r,i}=0$ for $i<0$ and $D_{r,s}(0)=1$. The later binomial determinants use a separate lower-case symbol $d$.
For the auxiliary problem, prescribe simple evaluations of polynomials at distinct nonzero nodes, together with the first $a$ Taylor coefficients at zero. Suppose the simple evaluations cancel in two groups of columns. How does one recover the signed leading coefficient of the entire determinant when the surviving columns within a group have the same large-frequency limit?

Two elementary steps will be used repeatedly. First, suppose the $a=u+v$ surviving columns in the jet space $\mathbb R[x]/(x^a)$ are $x^iP$ ($0\le i<u$) followed by $x^iQ$ ($0\le i<v$), with $P(0)\ne0$. Their ordered coefficient determinant equals
\begin{equation}\label{eq:two}
 P(0)^a\det\big([x^{u+i-j}]Q/P\big)_{0\le i,j<v}.
\end{equation}
Indeed multiplication by $P$ is triangular with diagonal $P(0)$. After division by $P$, the first $u$ columns are standard basis vectors. Second, if
$[x^h](Q/P)=H^{2h-1}(b_h+O(H^{-1}))$ at all required indices, the leading minor in \eqref{eq:two} has scale $H^{2uv-v}$ and signed coefficient $(-1)^{\binom v2}\det(b_{u-v+1+i+j})_{i,j<v}$. Every permutation contributes the same total index $uv$. If $b_h$ are moments of one sign of a positive measure with infinite support on $(0,\infty)$, this last determinant is nonzero whenever the displayed moments exist and $u\ge v$. The corresponding quadratic form is the integral of a squared polynomial against a positive power of the variable times that measure. These steps identify a leading coefficient even when the termwise approximation of the original evaluation determinant has dependent columns.

\subsection*{Moment and evaluation formulas}
Set $S_{-1}(y)=0$, $S_0(y)=1$, $S_{h+1}(y)=yS_h(y)-S_{h-1}(y)$, and use the monic bases
\[
 E_h(x)=S_h(x-2),\qquad O_h(x)=S_h(x-2)+S_{h-1}(x-2).
\]
For either basis use the semi-infinite symmetric tridiagonal matrix $J$ with off-diagonal entries $1$ and diagonal entries $2$, except that $J_{00}=1$ for $O$. With $L(f)=(f(J))_{00}$, the recurrence gives $B_i(J)e_0=e_i$ and $L(B_i B_j)=\delta_{ij}$ for $B=E,O$. The resolvent components $G_b(z)=\sum_{u\ge0} L(x^u B_b)z^u$ satisfy
\[
 G_0=1+z(J_{00}G_0+G_1),\quad
 G_b=z(G_{b-1}+2G_b+G_{b+1})\ (b\ge1).
\]
Substitution and uniqueness coefficient by coefficient give $G_b=z^b C^{2b+2}$ for $E$, and $G_b=z^b C^{2b+1}$ for $O$.
Hence, for $a\ge0$, the original shifted moments $c_{r,u+s}$ equal $L(x^u q)$, where
\[
 \begin{array}{c|ccc}
 r&B& b&q(x)\\\hline
 2t&E&t-1&x^a E_b(x),\quad a=s+t-1\\
 2t+1&O&t&x^a O_b(x),\quad a=s+t .
 \end{array}
\]
The distinct simple roots of $B_b$ are $2+2\cos(\ell\pi/t)$ for $1\le\ell<t$ in the even case and $2+2\cos(2\ell\pi/(2t+1))$ for $1\le\ell\le t$ in the odd case. Indeed at $x=2+2\cos\theta$,
\[
 E_h=\sin((h+1)\theta)/\sin\theta,\qquad
 O_h=\sin((h+1/2)\theta)/\sin(\theta/2).
\]
Put $d=a+b$. Let $V(f)$ list the first $a$ Taylor coefficients in ascending degree, then the evaluations at the $b$ simple roots in a fixed order. For any $N\ge0$,
\begin{equation}\label{eq:coll}
 D_{r,s}(N)=(-1)^{Nd}\frac{\det(V(B_{N+i}))_{0\le i<d}}{\det(V(B_i))_{0\le i<d}} .
\end{equation}
The denominator is nonzero by interpolation. To prove the identity, apply the $N$ functionals $L(B_i f)$ ($0\le i<N$) followed by $V$ on polynomials of degree $<N+d$. In the monic basis $B_0,\ldots,B_{N+d-1}$ expansion gives the numerator. In the monic basis $1,\ldots,x^{d-1},q,\ldots,x^{N-1}q$ block interchange gives $(-1)^{Nd}$ times the denominator times the original moment determinant. For a general monic orthogonal basis with nonunit squared norms the quotient in \eqref{eq:coll} must be multiplied by the product of the first $N$ norms, by the same block calculation.

The polynomials in both bases can be continued to all integer indices by recurrence. Along $N=tn+j$ in the even case, or $N=(2t+1)n+j$ in the odd case, each expression on the right of \eqref{eq:coll} is a polynomial in $n$ times $(-1)^{n\binom t2}$ or $(-1)^{tn}$ respectively. To verify this, the coefficients
\[
 [x^i]E_h=(-1)^{h-i}\binom{h+i+1}{2i+1},\qquad
 [x^i]O_h=(-1)^{h-i}\binom{h+i}{2i}
\]
follow from the generating function of $S_h$ and Pascal's identity; the backwards recurrence gives the same formulas at negative indices when generalized binomials are used. Each jet row thus has a common periodic sign. Increasing $n$ by one multiplies a simple row in the even case by $(-1)^\ell$ and leaves it fixed in the odd case. Including the sign in \eqref{eq:coll} gives exactly the stated normalizations. Since the normalized values are integers for every sufficiently large $n$, interpolation gives rational coefficients.

\subsection*{The three limits used to extract the coefficient}
All expansions here can be taken uniformly on a sufficiently small fixed circle about $z=0$. Thus estimates of the stated order hold for each fixed Taylor coefficient by Cauchy's formula. Put $x=z/H^2=2-2\cos\phi$. Then $\phi=\sqrt z/H+O(H^{-3})$, in the analytic variable $\sqrt z$. For integral $H>0$ put $P_H=E_{H-1}$, and for fixed nonzero integral $\delta$ let $R_H=P_{H+\delta}/P_H$. The sine addition formula gives
\[
 \frac{R_H(z/H^2)}{R_H(0)}
 =\frac{H}{H+\delta}\{\cos(\delta\phi)+\cot(H\phi)\sin(\delta\phi)\}
 =1+\frac{\delta}{H}(\sqrt z\cot\sqrt z-1)+O(H^{-2}).
\]
For a half-integer $M>0$ and integer $K=M+O(1)$, set
$U=\cos(M\phi)/\cos(\phi/2)$ and $W=\sin(K\phi)\tan(\phi/2)$.
At $x=z/M^2$ addition and Taylor expansion yield
\[
 W/U=\frac{\sqrt z\tan\sqrt z}{2M}+O(M^{-2}),\qquad
 U/(W/x)=\frac{2\sqrt z\cot\sqrt z}{M}+O(M^{-2}).
\]
The apparent singularities at the origin are removable. These limits, with the signed basis changes below, will supply both the magnitude and the sign.

\section{The five odd-power identities of Cigler}
Let $p=2t+1$ with $t\ge1$ and write $H(N)=D_{p,0}(N)$.
We prove the identities in Cigler, \emph{Catalan numbers, Hankel determinants
and Fibonacci polynomials}, arXiv:1801.05608v3, Conjecture~7.2, equation (7.6):
\begin{align*}
 H(pn)=H(pn+1)&=(-1)^{tn},&
 H(pn+t+1)&=0,\\
 H(pn+t)=-H(pn+t+2)&=(-1)^{tn+\binom t2}\big(p(n+1)\big)^{t-1},\\
 H(pn-1)+H(pn+2)&=(-1)^{tn+1}p(t-1).
\end{align*}
The first four evaluations hold for $n\ge0$; the final identity between ordinary determinants holds for $n\ge1$. The extension to $n=0$ used in its proof is defined below. Both powers in the final sum are $p$ in Cigler's original. Wang and Xin's Conjecture~7 printed $2t$ in its last determinant. The latter version fails for $t=n=1$, since $D_{3,0}(2)=0$, whereas $D_{2,0}(5)=1$.

Wang and Xin, arXiv:1811.00248, Theorem~11, investigated the conjectures up to $r=31$. The $H(pn+t)$ evaluation also follows by setting the shift to zero in Cigler's later Conjecture~11, as proved in the trureturing record \emph{Shifted Hankel Determinants of Odd Catalan Powers}. In the earlier Stieltjes-tableau notation of Cigler's \emph{Some remarks and conjectures about Hankel determinants}, $a(u,t,2,-1)=[z^{u-t}]C(z)^p$, so $d_s(N,t,2,-1)=D_{p,s-t}(N)$. The results there for $s=0,1,2$ meet the unshifted case only when $s=t$. Fulmek, arXiv:2402.19127, establishes identities relating backward and forward shifts. We give a direct argument for all five displayed evaluations.

In \eqref{eq:coll} the multiplier is $x^t O_t$, so $H(N)$ is the ratio of the two evaluation determinants with $t$ jets and $t$ simple nodes; the external sign is one. The same quotient gives $H(-1)=0$ because $O_{-1}=-O_0$. The denominator equals $v=\det(O_i(r_\ell))_{\ell,\,0\le i<t}$: use the ordered basis $1,\ldots,x^{t-1},x^tO_0,\ldots,x^tO_{t-1}$ and $\prod_\ell r_\ell=1$, the latter following from $O_t(0)=(-1)^t$.

Index each column by its half-integral frequency $u=N+s+1/2$. At the simple roots its values are $\sin(2\pi l u/p)/\sin(\pi l/p)$; at $x=2-2\cos\phi$ its jet is obtained from
\[
 O_{u-1/2}(x)=(-1)^{u-1/2}\cos(u\phi)/\cos(\phi/2).
\]
Put $M=(n+1/2)p$, $K=(n+1)p$, $U_M=\cos(M\phi)/\cos(\phi/2)$ and
$W_K=\sin(K\phi)\tan(\phi/2)$. Then $U_M(0)=1$ and $[x]W_K=K/2$.
At $M$ the singleton and pair sums at $M\pm i$, $1\le i<A$, give $A$ columns with zero simple-root entries and jets equal to $U_M$ times polynomials of successive degrees $0,\ldots,A-1$. Indeed the factors are signed multiples of $1,\cos\phi,\ldots,\cos((A-1)\phi)$, with a common sign changing by $-1$ upon increasing $n$. At the integral center $K$ the pair sum at $K\pm(i+1/2)$ has zero simple-root entries and equals
\[
 -2(-1)^{K+i}W_K\left(1+2\sum_{j=1}^i\cos(j\phi)\right).
\]
Its polynomial factor has degree $i$ and leading coefficient $-2(-1)^K$ in $x$. These column operations keep one representative of each pair.

For $N=pn,pn+1$ use the group $A=t$ at $M$. Multiplication by $U_M$ has determinant one on the jets modulo $x^t$. The complementary simple-root block stays fixed upon increasing $n$, since each frequency increases by $p$. The only variation is the sign $(-1)^t$ per step from the $t$ kernel columns. The initial values $H(0)=H(1)=1$ prove the first identities. For $N=pn+t+1$ the frequencies are exactly $K\pm(i+1/2)$, $0\le i<t$. Their $t$ pair sums have zero simple-root rows and jets divisible by $x$, a space of dimension $t-1$. This proves the vanishing.

For $N=pn+t$ the first kernel column has frequency $M$. The other $t-1$ kernel columns are the pair sums at $K\pm(i+1/2)$, $0\le i<t-1$, formed by adding the left column to the right column. The $t$ remaining original columns have indices $N+1,\ldots,N+t$. Moving them behind the $t-1$ sums gives the sign $(-1)^{t(t-1)}=1$. Expand on the simple rows. At $n=0$ their block is the negative of the reversed matrix of $O_0,\ldots,O_{t-1}$ at the nodes, by $\sin((p-u)\theta_l)=-\sin(u\theta_l)$, $\theta_l=2\pi l/p$. It supplies the sign $(-1)^{t+\binom t2}$ against the denominator, and stays fixed for all $n$. Expand the jet determinant on its constant row: only the singleton survives, with value $(-1)^{M-1/2}$. For the remaining columns multiplication by $W_K/x$ contributes $(K/2)^{t-1}$, and the triangular polynomial factors contribute $(-2(-1)^K)^{t-1}$. At $n=0$ this jet determinant is $(-1)^t p^{t-1}$ and its sign changes by $(-1)^t$ per step in $n$. This proves the formula for $H(pn+t)$.

Reflection of all these frequencies about $K$ gives those for $N=pn+t+2$ in reverse order. Reversing $2t$ columns and negating $t$ simple rows multiply the determinant by $(-1)^{\binom{2t}{2}+t}=1$ and restore the same simple-row array after reflection. The $t-1$ pair sums at $K$ remain unchanged. The singleton now has frequency $M+p$ and opposite constant term. Expansion on the same blocks proves $H(pn+t+2)=-H(pn+t)$.

Finally, for $t=1$ both indices in the sum are $2$ modulo $3$ and the vanishing just proved applies. For $t\ge2$ start with the interval for $N=pn+2$. Its $t$ kernel columns consist of the $A=t-1$ group at $M$ and the single pair sum at $K\pm1/2$. Reflect the whole interval by $u\mapsto 2M-u$. This is the reverse of the interval for $pn-1$. Once again column reversal and simple-row negation have total sign $+1$. After those operations both arrays have identical simple blocks and the same $t-1$ kernel columns centered at $M$. The final pair for the reflected array is centered at $K'=np$.
It follows by multilinearity and expansion on the simple rows that the sum of the two determinants uses a single last jet column proportional to $(-1)^K(W_K-W_{K'})$, with proportionality constant independent of $n$; here $K-K'=p$ is odd. The $t-1$ common kernel columns have the form $U_M$ times the polynomial basis of degrees $0,\ldots,t-2$, each with factor $(-1)^n$ relative to $n=0$. Their simple complement is independent of $n$ by periodicity. Divide every jet column by the unit $U_M$; this leaves the jet determinant unchanged. The last column, apart from the sign $(-1)^K$, is independent of $n$ because
\[
 (W_K-W_{K'})/U_M
 =2\sin(p\phi/2)\sin(\phi/2).
\]
Thus the sum is $(-1)^{tn}$ times its value at $n=0$. The bilateral value $H(-1)=0$ and the first three coefficients of $C^p$ give
$H(2)=p(p+3)/2-p^2=-p(t-1)$. This determines the claimed constant and proves the last identity.

In fact, the same reduction yields the entire residue polynomial at $j=2$. For $t\ge2$, with $x=2-2\cos\phi$ and $M=(n+1/2)p$, it is
\[
 (-1)^{tn}H(pn+2)
 =2(-1)^{t-1}[x^{t-1}]\sin(\phi/2)
 \left(\cos(p\phi/2)\tan(M\phi)+\sin(p\phi/2)\right).
\]
Indeed expansion on the $t-1$ monomials leaves $[x^{t-1}]W_K/U_M$. To retain the scalar, reorder the kernel columns as the increasing offsets from $M$, then the single pair at $K$, then the simple frequencies $M+1,\ldots,M+t$. The permutation sign is $(-1)^{\binom{t-1}{2}+t}$, the simple block gives $(-1)^{t+\binom t2}$, the half-integral-center factors give $(-1)^{(M-1/2)(t-1)}$, and the last pair gives $-2(-1)^K$. Combined with $(-1)^{tn}$ their product is $2(-1)^{t-1}$. The contribution of the second sine is $-p(t-1)/2$, since $2\sin(p\phi/2)\sin(\phi/2)=\cos(t\phi)-\cos((t+1)\phi)$ and $[x^j]\cos(q\phi)=(-1)^j q\binom{q+j}{2j}/(q+j)$ for integers $q\ge1$, $j\ge0$; the latter follows from the Chebyshev recurrence. The other contribution is an odd polynomial in $n+1/2$. Thus adding $p(t-1)/2$ to the normalized residue polynomial makes its odd symmetry about $-1/2$ explicit.

\section{Cigler's even-power determinant identities}
For $k\ge1$ put $H_e(N)=D_{2k,0}(N)$. Cigler, arXiv:1801.05608v3,
Conjecture~7.5, equation (7.11), states two identities for $k>0$.
The first holds for every $k\ge1$ and $n\ge0$:
\[
 H_e(kn)=H_e(kn+1)=(-1)^{n\binom{k}{2}}(n+1)^{k-1}.
\]
The second holds for every $k\ge2$ and $n\ge1$:
\begin{equation}
 H_e(2kn-1)+H_e(2kn+2)=-k(2k-3)(2n+1)^{k-1}.\tag{C7.5}
\end{equation}
The stated $k=1$ endpoint of the second equation fails: $C(z)^2=(C(z)-1)/z$, and the Hankel determinants of the shifted Catalan numbers are all $1$. Thus the left side is $2$ and its proposed value is $1$. In particular $k=1,n=1$ suffices. The endpoint discrepancy has previously been recorded \cite{TT}. Wang and Xin treated small powers by continued fractions.

For completeness let $E_j(x)=S_j(x-2)$ and index its column by the frequency $u=j+1$, writing $P_u=E_{u-1}$. The even Christoffel quotient has $a=k-1$ jets at zero and $a$ simple nodes $x_l=2+2\cos(l\pi/k)$, $1\le l<k$. Its multiplier is $x^a E_{k-1}$ and its exterior determinant sign is one because the total degree $2a$ is even. The constant denominator is nonzero. At the nodes the frequency column is $\sin(ul\pi/k)/\sin(l\pi/k)$, and with $x=2-2\cos\phi$ its jet comes from
\[
 P_u(x)=(-1)^{u-1}\sin(u\phi)/\sin\phi,\qquad P_u(0)=(-1)^{u-1}u.
\]
These expressions follow from the same tridiagonal construction, now with diagonal $2$ even at index zero. All its off-diagonal entries equal $1$, hence its monic orthogonal polynomials $E_j$ have norms one. In particular $k=1$, where there are no rows in the quotient, gives $H_e(N)=1$ for all $N$.

Take $k\ge2$. For $N=kn,kn+1$ the frequency range from $N+1$ to $N+2k-2$ contains $M=k(n+1)$ and $M\pm i$ for $1\le i\le k-2$. Make the singleton $P_M$ and the pair sums at those frequencies the $k-1$ kernel columns. They vanish at every simple node and are $P_M$ times polynomials of successive degrees $0,\ldots,k-2$ with fixed nonzero leading coefficients: the factors for pairs are $2\cos(i(\pi-\phi))$. Compared with $n=0$, the simple-row block has row signs $(-1)^{nl}$, while the jet block contributes the relative factor
$\{P_{k(n+1)}(0)/P_k(0)\}^{k-1}$.
As $k(k-1)$ is even, the latter has positive sign. Using $H_e(0)=H_e(1)=1$ gives the asserted first equality.

For the second equation use $M=(2n+1)k$, $K=(2n+2)k$, and $K'=2nk$. The frequency interval for $N=2kn+2$ is $[2kn+3,2kn+2k]$. Its reflection about $M$ gives the reverse of the interval for $N=2kn-1$. Reversal of the $2k-2$ columns has sign $(-1)^{k-1}$. Reflection negates all $k-1$ simple rows, since $2M$ is divisible by $2k$. Those signs cancel: after reversing columns and negating simple rows, the second determinant has the same simple entries, in the same positions, as the first. For $k\ge3$ form $k-2$ common kernel columns at $M$: the singleton and the sums at $M\pm i$ for $1\le i<k-2$. For $k=2$ this group is empty. In either case the last singleton kernel column is $P_K$ in the first determinant and $P_{K'}$ in the second, and the reflected common columns are unchanged. Keeping one original representative of each pair gives the identical complementary simple block for both arrays. Its entries are independent of $n$ by periodicity. Expanding on the simple rows and adding the determinants therefore replaces the last kernel column by $P_K+P_{K'}$. Since $K$ and $K'$ are even,
\[
 P_K+P_{K'}=2(-1)^{K-M} P_M\cos(k\phi).
\]
Thus every kernel jet column is $P_M$ times a fixed polynomial in $x$: multiplication in $\mathbb R[x]/(x^{k-1})$ contributes exactly $P_M(0)^{k-1}$. The ratio of this scalar to its value at $n=0$ is $(2n+1)^{k-1}$. At $n=0$ the continued numerator for $H_e(-1)$ contains the zero column $P_0$, while $H_e(2)=k(2k+3)-(2k)^2=-k(2k-3)$. This proves (C7.5) for all $k\ge2$.

\section{Transport of signed leading coefficients for even powers}
The next computation retains the normalization constants in the evaluation determinant. It proves the two Bernoulli formulas of Cigler's Conjecture~8 and the entire shift law in his Conjecture~9 (arXiv:2308.07642v2, equations (22)--(24) and the paragraph immediately following them). For $h\ge1$ define
\[
 c_h=2\sum_{s=1}^\infty(\pi s)^{-2h},\qquad
 b_h=2\sum_{s=1}^\infty(\pi(s-1/2))^{-2h},\qquad
 \mathcal H_v(f,w)=\det(f_{w+i+j})_{0\le i,j<v}.
\]
In particular $-c_h=[z^h](\sqrt z\cot\sqrt z)$ and $b_h=[z^h](\sqrt z\tan\sqrt z)$, by the products for sine and cosine. The empty determinant equals one. For $w\ge1$ and $f=b,c$ every $\mathcal H_v(f,w)$ is positive because the quadratic form for $f=c$ equals $2\sum_{s\ge1}(\pi s)^{-2w}P((\pi s)^{-2})^2$ for a polynomial $P$ of degree $<v$, and $s$ is replaced by $s-1/2$ for $f=b$.

Fix $k\ge3$, $2\le j\le k-1$, and $m\ge-2$. Put
\[
 t=k+m,\quad A=k-j,\quad B=j-1,\quad a=A+B=k-1,\quad
 U=\max(A,B),\quad V=\min(A,B).
\]
For the normalized determinant
\[
 f_m(n)=(-1)^{n\binom t2}D_{2t,-m}(tn+j+m)
\]
the degree is $(2V+1)U$, and its signed leading coefficient is
\[
 \operatorname{lc} f_m=(-1)^{e(t)}t^{2AB}\mathcal H_V(c,U-V+1).
 \tag{E1}
\]
The sign exponent, used only modulo two, has the following explicit formula. Set
\[
 e_0(t)=(t-A)(a+t-1)+(t-1)B+\binom A2+\binom t2.
\]
Then
\[
 e(t)=e_0(t)+
 \begin{cases}
 (t-1)a+t B+B+\binom B2,&A\ge B,\\
 AB+(2t-1)a+t A+\binom A2,&A<B.
 \end{cases}
 \tag{E2}
\]
The case $t=1$ can occur at $(k,j,m)=(3,2,-2)$; all simple-root determinants below are then empty, with the same formulas.

Here the moment multiplier is $x^aE_{t-1}$, so the simple-root count is $t-1$ and the total column count is $d=a+t-1$. Set $P_u=E_{u-1}$, $M=t(n+1)$, $K=t(n+2)$, but now the order is $N=tn+t-A$. The column frequencies range from $M-A+1$ to $K+B-1$. The bounds on $j,m$ imply $1\le A,B\le t$. The three consecutive groups consist of the $A$ columns ending at $M$, the $t-1$ middle columns $M+1,\ldots,K-1$, and the $B$ columns starting at $K$. Replace each strictly exterior frequency $M-i$ by the sum of its column and the middle column $M+i$; replace $K+i$ by its sum with $K-i$. Every required middle column exists even when $m=-2$. Retain the two singleton columns at $M$ and $K$. Ordering the resulting columns as the $M$ cluster by increasing $i$, then the $K$ cluster by increasing $i$, then the middle columns gives sign $(-1)^{\binom A2+(t-1)B}$. The sine identity $P_{H+i}+P_{H-i}=P_H\,(S_i(x-2)-S_{i-2}(x-2))$ for $i\ge1$ has a monic factor of degree $i$. These pair identities make both clusters vanish on the simple-root rows. Since their polynomial factors are monic by degree, their zero-jet columns may be taken as $x^iP_M$, $0\le i<A$, and $x^iP_K$, $0\le i<B$, without changing the determinant.

Use the jet-first row order in numerator and denominator. Write $v$ for the determinant of evaluation of $1,x,\ldots,x^{t-2}$ at the simple nodes, with the same row order. Monicity gives the same determinant for $E_0,\ldots,E_{t-2}$. By taking monomials in the denominator, its value is $v\prod_{\ell}\lambda_\ell^a=v t^a$, since $E_{t-1}(0)=(-1)^{t-1}t$. In the middle block, evaluation of $P_{M+i}$ at the $\ell$-th root equals $(-1)^{(n+1)\ell}$ times evaluation of $E_{i-1}$, hence its determinant is $(-1)^{(n+1)\binom t2}v$. This proves that the exact remaining scalar is $t^{-a}$ times the signs already recorded and the Christoffel sign $(-1)^{Nd}$.

If $A\ge B$ take $H=M,\delta=t$. Extract $P_H(0)^a$ from the jet determinant using \eqref{eq:two}, and take out $R_H(0)^B$ from the $B$-row minor, where $R_H=P_{H+\delta}/P_H$. The signs of these two factors are $(-1)^{(M-1)a+tB}$. Positive-index entries of $R_H/R_H(0)$, after division by the appropriate $H^{2h-1}$, tend to $-\delta c_h$, so the minor contributes sign $(-1)^{B+\binom B2}$: the binomial exponent comes from reversing its columns. For $A<B$, swap the two clusters, giving sign $(-1)^{AB}$. Now $H=K,\delta=-t$; the extracted signs are $(-1)^{(K-1)a+tA}$, and the leading minor contributes $(-1)^{\binom A2}$ because $-\delta>0$. In both cases the absolute leading scale before changing variable from $H$ to $n$ is
\[
 t^{-a}t^V H^{a+2AB-V}\mathcal H_V(c,U-V+1).
\]
Since $H=tn+O(1)$ this is precisely the absolute value in (E1). These are genuine leading terms because the moment determinants are strictly positive. In the signs involving $n$, the row sign, $(-1)^{Nd}$, the factors $P_H(0)^a$ and the normalization together have even total coefficient of $n$; equivalently $t(d+a)=t(t-1+2a)$ is even. The constant sign is (E2).

Holding $k,j$ fixed, direct subtraction in (E2) gives
\[
 e(t+1)-e(t)\equiv t+A\equiv m+j\pmod2.
\]
Indeed the first two terms of $e_0$ have changes of parity $a-A=B$ and $B$, its binomial term in $t$ changes by $t$, and either branch contributes $A$. Iterating in both directions, and using the ordinary polynomial definition of $\binom m2$ for negative integers, gives
\[
 \frac{\operatorname{lc}f_m}{\operatorname{lc}f_0}
 =\left(\frac{k+m}{k}\right)^{2(j-1)(k-j)}
 \begin{cases}(-1)^{\binom m2},&j\ \text{even},\\
 (-1)^{\binom{m+1}{2}},&j\ \text{odd}.
 \end{cases}
\]
This proves Conjecture~9 without an assumed value or sign for the reference coefficient. It also determines the constants in Conjecture~7 for every one of its defined residue indices:
\[
 A_{2k,j}=(-1)^{e(k)}\mathcal H_{\min(k-j,j-1)}
 \bigl(c,|k-2j+1|+1\bigr),\qquad 1\le j\le k-1 .
\]
Here the displayed definition (16) in Cigler assigns residue indices $0,\ldots,k-1$; Conjecture~7 asks for $j\ge1$ among those indices. At $j=1$, set $m=0$, $A=k-1$ and $B=0$ in the same column proof: the right cluster is absent, the root-node normalizer is $k^a$, and the single jet generator gives $P_M(0)^a$. The resulting polynomial is $(n+1)^{k-1}$, giving $A_{2k,1}=1$. For the translated choice $j=k$ the left cluster is absent and the same computation gives $(-1)^{\binom k2}(n+2)^{k-1}$. Thus the residue zero itself has $(n+1)^{k-1}$ and leading coefficient one. For $k=1$ the determinant is identically one. This explicit evaluation supplies information beyond the existence of a freely named scalar in Conjecture~7.

For Conjecture~8 take $m=0$ and $j=2$ or $k-1$. Then $V=1$, $U=k-2$ and the unsigned moment determinant is $c_{k-2}$. Formula (E2) gives $e(k)\equiv k$ for $j=2$ and $e(k)\equiv k+\binom{k+2}{2}$ for $j=k-1$ (the latter also holds for $k=3$). The identity
\[
 c_h=\frac{2^{2h}|B_{2h}|}{(2h)!},\qquad
 \operatorname{sgn} B_{2h}=(-1)^{h+1}\quad(h\ge1)
\]
therefore gives for the two sign-normalized sequences in equations (22) and (23) exactly the common leading coefficient $-(2k)^{2k-4}B_{2k-4}/(2k-4)!$, at degree $3(k-2)$.

\section{Signed odd coefficients and the Bernoulli sign in equation (26)}
Let $p=2t+1$, $1\le j\le t$, $A=t-j+1$, $B=j-1$ and $D=B(2A-1)$. Write $L_j$ for the leading coefficient of $(-1)^{tn}D_{p,0}(pn+j)$, using $j=p$ as a translate of residue zero. The full coefficient, including its sign, is
\[
 \frac{L_j}{p^D}=
 \begin{cases}
 (-1)^{AB}\mathcal H_B(b,A-B+1),& A\ge B,\\[2pt]
 (-1)^{\binom{B+1}{2}+\binom A2}\mathcal H_{A-1}(c,B-A+2),& B\ge A-1 .
 \end{cases}
 \tag{O1}
\]
Either expression is valid in their overlap. The mirror relation is
\[
 L_{p+1-j}=(-1)^A L_j.\tag{O2}
\]
Together with the zero middle residue, these formulas prove the leading-coefficient assertion and mirror signs of Cigler's Conjecture~13. They also identify the precise sign in each Bernoulli formula in Conjecture~14.

For normalization use the same orthogonal determinant as in the proof of Conjecture~7.2. The denominator's product of all simple roots is $1$ since $O_t(0)=(-1)^t$. Its determinant is hence just $v$, the Vandermonde at the $t$ simple nodes. As before use $U(x)=\cos(M\phi)/\cos(\phi/2)$ and $V(x)=\sin(K\phi)\tan(\phi/2)$ for $x=2-2\cos\phi$. In this section $U,V$ denote these polynomials, rather than multiplicities. Conversion of the $A$ half-integral-center columns at $M$ to $x^iU$ contributes exactly the sign $(-1)^{(M-1/2)A}$. Indeed the singleton has that one-column sign $(-1)^{M-1/2}$, and for offset $i>0$ the factor multiplying $U$ is $(-1)^{M-1/2+i}2\cos(i\phi)$, whose $x^i$ coefficient equals $(-1)^{M-1/2}$. At the integer center $K$, the pair at displacement $i+1/2$ gives
\[
 2(-1)^{K+i+1}V\,\frac{\sin((i+1/2)\phi)}{\sin(\phi/2)} .
\]
The last quotient has degree $i$ and leading $x$ coefficient $(-1)^i$. The net conversion factor for the $B$ columns is therefore $2^B(-1)^{(K+1)B}$. These factors account for all scales from the kernel columns.

At the lower residue the full half-integral frequency interval runs from $M-A+1$ to $K+B-1/2$, and at the mirror residue from $K-B+1/2$ to $M+A-1$. At the first center use its singleton and the pairs at displacements $1,\ldots,A-1$; at the integer center use displacements $1/2,\ldots,B-1/2$. In each interval retain the intervening $t$ original columns and add the matching retained column to each exterior member of a pair. For the lower residue $j$, use $M=(n+1/2)p$ and $K=(n+1)p$. Order the columns by the ascending offsets from $M$, then the ascending offsets from $K$, then the central frequencies $M+1,\ldots,M+t$. This permutation contributes sign $(-1)^{\binom A2+tB}$. At $\theta=2\pi\ell/p$, the central frequency $M+i$ gives minus the evaluation at frequency $p/2-i$. These last frequencies are $t-1/2,\ldots,1/2$ in reverse order. Consequently the simple-root block equals $(-1)^{t+\binom t2}v$. In the mirror residue $p+1-j$, the left center is $K=(n+1)p$ and the right center is $M=(n+3/2)p$. Ordering the columns in the same order of generators ($U$ then $V$) and moving the middle frequencies last gives sign $(-1)^{A(B+t)+\binom B2}$. The simple-root frequencies are now $K+1/2,\ldots,K+t-1/2$, giving the determinant $v$ itself. The Christoffel degree $2t$ is even. Combining these permutations and conversions with the normalization $(-1)^{tn}$ leaves a common absolute scalar $2^B$, multiplying the $U,V$ jet determinant, with residual sign exponents respectively
\[
 e_{\rm low}=\binom B2+AB,\qquad
 e_{\rm high}=\binom B2+AB+A .
\]
For example the first reduction uses $M-1/2\equiv n+t$, $K+1\equiv n$ and $\binom t2=\binom A2+\binom B2+AB$ modulo two. The two sets of frequencies differ in the displacement of the centers, but $K/M\to1$ in both. Thus (O2) follows once the common asymptotic normalization is calculated.

For $A\ge B$, $U(0)=1$ and the $B$-row minor contains Taylor coefficients of $V/U$. The limit after $x=z/M^2$ is $(2M)^{-1}\sqrt z\tan\sqrt z+O(M^{-2})$. It cancels the exterior $2^B$ exactly in the leading coefficient. Column reversal contributes $(-1)^{\binom B2}$ and all moment entries are $b_h$, proving the first line of (O1). If $B\ge A-1$, delete the constant row and the first $U$ column, divide the other columns by $x$, and move all $B$ generators $V/x$ before the $A-1$ remaining generators $U$. This costs sign $(-1)^{B(A-1)}$. Dividing by $V/x$ extracts $(K/2)^{t-1}$ and the scaled ratio $U/(V/x)$ is $(2/M)\sqrt z\cot\sqrt z+O(M^{-2})$. The powers of $2$ cancel because $B+(A-1)=t-1$. The sign from the cotangent coefficients and the reversal is $(-1)^{A-1+\binom{A-1}{2}}$. Combining with $e_{\rm low}$ gives the second sign in (O1), while $M=pn+O(1)$ gives $p^D$ in both cases. The estimates are uniform on a small complex circle and the moments are strictly positive, by the moment criterion. At the middle residue $j=t+1$ all $t$ pair sums are centered at the integer $K$; their $V$ factors vanish at $x=0$. The simple-root block still has full rank, while the $t$ kernel jet columns have zero constant row. Hence the middle determinant vanishes for all $n$.

For $t\ge2$ specialize $j=2$, giving $B=1$, $A=t-1$. Equations (O1)--(O2) yield
\[
 (-1)^{t-1}L_2=L_{2t}
 =p^{2t-3}\frac{2^{2t-2}(2^{2t-2}-1)|B_{2t-2}|}{(2t-2)!}.
 \tag{O3}
\]
The absolute value is necessary under the signed Bernoulli convention used by Cigler. In particular his equation (26), which uses $B_{2t-2}$ without that absolute value, gives $-343/3$ at $t=3$, whereas (O3) gives $+343/3$. This is a disproof of the literal Conjecture~14; (O3) states the corrected uniform signed formula separately.

The other equation in that conjecture, (27), holds in its natural range $t\ge2$. For $t\ge3$ take $j=t-1$, $A=2$, $B=t-2$. The second line of (O1) then gives one cotangent moment at index $t-2$. Substitution yields
\[
 L_{t-1}=(-1)^{\binom{t-2}{2}}
 p^{3t-6}\frac{2^{2t-4}B_{2t-4}}{(2t-4)!}.
\]
The same formula for $t=2$ reads $L_1=1$, since $B_0=1$, and follows from the empty tangent determinant in (O1).

\section{The optimal common denominator of the normalized coefficients}
Define $\Phi_k=\prod_{i=1}^{k-1}(2i-1)!!$ for $k\ge1$, with empty product one. In the remark after Conjecture~7, Cigler exhibits integral multiples of $1/\Phi_k$ in the first rows for the even family. For every $k\ge2$, the least positive integer clearing the denominators of all $A_{2k,j}$, $1\le j<k$, is exactly $\Phi_k$. For odd powers $2t+1$ with $t\ge1$, the least positive integer clearing all nonzero normalized coefficients $L_j/(2t+1)^{D}$ is exactly $\Phi_t$, where the degree $D$ is that of the corresponding residue.

The key identity applies simultaneously to arbitrary shifts of the positive moment determinants above. Fix $u\ge v\ge0$, $a=u+v\ge1$, and let $\mathcal E$ be the increasing rearrangement of the $a$ distinct integers
\[
 0,2,\ldots,2u-2,\qquad 1,3,\ldots,2v-1.
\]
Write $(x)_{\underline e}=x(x-1)\cdots(x-e+1)$ and $\Delta(X)=\prod_{i<j}(X_j-X_i)$. Then
\[
\begin{aligned}
 \Phi_{a+1}\mathcal H_v(c,u-v+1)
 &=\left|\frac{\det((X_i)_{\underline e})_{0\le i<a,\ e\in\mathcal E}}{\Delta(X)}\right|_{X_i=2i+1},\\
 \Phi_a\mathcal H_v(b,u-v+1)
 &=\left|\frac{\det((X_i)_{\underline e})_{0\le i<a,\ e\in\mathcal E}}{\Delta(X)}\right|_{X_i=2i}.
\end{aligned}
\]
Each quotient is the evaluation of a symmetric polynomial with integer coefficients and has strictly positive absolute value.

To see this, take the determinant of the coefficients of $z^0,\ldots,z^{a-1}$ in the columns $z^i F$ for $0\le i<u$ and $z^j G$ for $0\le j<v$. Use
\[
 (F,G)=\begin{cases}
 (\sin\sqrt z/\sqrt z,\,\cos\sqrt z)&\text{for }c,\\
 (\cos\sqrt z,\,\sqrt z\sin\sqrt z)&\text{for }b.
 \end{cases}
\]
Since $F(0)=1$, division by $F$ preserves the determinant, and expansion on the first $u$ columns followed by reversal gives absolute value $\mathcal H_v(c,u-v+1)$ or $\mathcal H_v(b,u-v+1)$. In the first case multiply row $r$ by $(2r+1)!$. After removing alternating row and column signs the entries become the falling factorials at $X_r=2r+1$ of degrees $2i$ in the first group and $2j+1$ in the second. In the second case multiply row $r$ by $(2r)!$ and the same degrees occur at $X_r=2r$. If a coefficient index is below the support of its series, the resulting falling factorial is zero as required. Sorting the degrees affects only a sign.

For indeterminate $X$, the determinant vanishes identically when $X_j=X_i$. Each factor $X_j-X_i$ therefore divides it in $\mathbb Z[X_0,\ldots,X_{a-1}]$, and these pairwise nonassociate prime elements have product $\Delta(X)$. On either arithmetic progression its value is $2^{\binom a2}\prod_{r=0}^{a-1}r!$. Dividing the product of the row factors by this Vandermonde gives $\Phi_{a+1}$ and $\Phi_a$, respectively, since $(2r+1)!/(2^r r!)=(2r+1)!!$ and $(2r)!/(2^r r!)=(2r-1)!!$. Positivity follows from the moment representation.

When $u-v\in\{0,1\}$, the degree set $\mathcal E$ is exactly $0,\ldots,a-1$. The polynomials are monic, so the determinant equals the Vandermonde and the quotient equals one. For even powers the indices $u,v$ in (E1) have sum $k-1$. The balanced choice $j=1+\lfloor(k-1)/2\rfloor$ gives $u-v\le1$ (for $k=2$ use $j=1$). Thus a coefficient of absolute value $1/\Phi_k$ occurs. For odd powers, the first branch of (O1) has $a=t$ in the tangent formula, while its second branch has $a=t-1$ in the cotangent formula. Both give the same denominator $\Phi_t$, with the case $a=0$ immediate. Taking $j=1+\lfloor t/2\rfloor$ in the tangent branch gives a balanced pair and attains $1/\Phi_t$ (at $t=1$ use $j=1$). The mirror signs preserve denominators. This proves the sharpness assertions.

Products for the two lowest shifts of closely related Bernoulli and Rayleigh Hankel determinants already occur in A.~Baricz and F.~\v{S}tampach, \emph{The Hurwitz-type theorem for the regular Coulomb wave function via Hankel determinants}, Linear Algebra Appl., DOI 10.1016/j.laa.2018.03.012, arXiv:1708.07729. Here the alternating polynomial applies to every $u-v\ge0$ and identifies the exact common denominators of the Catalan leading-coefficient families.

\section{Wang and Xin's binomial determinant conjecture}
Define
\[
 d_r(N)=\det\left(\binom{2i+2j+r}{i+j}\right)_{0\le i,j<N},
 \qquad d_r(0)=1.
\]
These are the Hankel determinants of $G(x,r)$ in Wang and Xin, Conjecture~22 of \cite{WX} (preprint numbering). We establish both parities of that conjecture.
For $r=2t+1$, $1\le j\le\lfloor(r+3)/4\rfloor$, the four determinants at orders
\[
 rn+j,\quad rn+t+j,\quad rn+t+2-j,\quad rn+r+1-j
\]
are polynomials in $n$ of exact degree $(2j-1)(t+1-j)$.
For $r=2t$, $1\le j\le t$, the determinants at orders $rn+j,rn+r+1-j$, multiplied by $(-1)^{tn}$, are polynomials of exact degree $(j-1)(r+1-2j)$. These assertions include $n=0$. Immediately after their conjecture the authors report that their method can verify the result through $r=26$.

\subsection*{Two positive functionals and the folded columns}
Use $x=2+2\cos\theta$ and let $L_0(f)$ be the average of $f(x)$ over $0\le\theta\le\pi$. Constant-term extraction with $w=e^{i\theta}$ gives
\[
 \binom{2h+r}{h}
 =\operatorname{CT}_w(2+w+w^{-1})^h(1+w)^r
 =L_0\big(x^h x^{r/2}\cos(r\theta/2)\big).
\]
For $r=2t$, $t\ge1$, use $L_0$ with monic basis $P_0=1$, $P_s=2\cos(s\theta)$ for $s\ge1$. Their squared norms are $1,2,2,\ldots$. The monic moment multiplier is $Q=x^t P_t$, accompanied by the factor $1/2$. Its zeros are $t$ zeros at the origin and the simple nodes
$\theta_l=(2l-1)\pi/r$, $1\le l\le t$. For $r=2t+1$ instead put $L_1(f)=L_0(xf/2)$ and use the monic basis
\[
 P_s=\cos((s+1/2)\theta)/\cos(\theta/2),\qquad s\ge0.
\]
Its squared norms are all $1$ by cosine orthogonality, and its monicity follows either from the recurrence or from the cosine addition formula. Here the monic moment multiplier is $Q=x^tP_t$ with no scalar factor, and its simple nodes again have angles $(2l-1)\pi/r$, $1\le l\le t$. For $t=0$ this immediately gives $d_1(N)=1$ for all $N$.

For $t\ge1$ in either parity let $V$ consist of the first $t$ Taylor coefficients at zero and the values at the $t$ simple nodes. The Christoffel determinant identity gives
\[
 d_r(N)=c_r\frac{\det\big(V(P_{N+s})\big)_{0\le s<2t}}{\det\big(V(P_s)\big)_{0\le s<2t}},\qquad N\ge1,
 \qquad c_{2t}=1/2,\quad c_{2t+1}=1.
\]
In the two-basis argument, the unperturbed block contributes the product of the first $N$ squared norms. In the even case its product $2^{N-1}$ combines with $(1/2)^N$ to give $c_{2t}$. All exterior signs are positive because $\deg Q=2t$.

Use frequency $u=N+s$ for even $r$ and $u=N+s+1/2$ for odd $r$. The value at each simple node is a fixed nonzero row factor times $\cos(u\theta_l)$, since $N\ge1$. At $x=2-2\cos\phi$ the jets are obtained from $2(-1)^u\cos(u\phi)$ in the even case, and from $(-1)^{u-1/2}\sin(u\phi)/\sin(\phi/2)$ in the odd case. Each fixed Taylor coefficient with the displayed sign removed is a polynomial in $u$, by the expansions of sine and cosine (in the odd case the removable denominator is linear in $\phi$ at zero). Shifting $n$ by $1$ multiplies every simple row by $-1$ since $r\theta_l$ is an odd multiple of $\pi$. The jet signs do not change for even $r$, and all change for odd $r$. Thus the sign-normalized determinants on each indicated residue are polynomials. Rationality of coefficients follows by interpolation at integers.

First take $N=rn+s$ with $1\le s\le t$ in the even case and $1\le s\le t+1$ in the odd case. Put
\[
 A=t+1-s,\quad B=s-1,\quad A+B=t,\qquad
 M=(n+1/2)r,\quad K=(n+1)r.
\]
At $M$ take the singleton and sums at $M\pm i$ for $1\le i<A$ when $A>0$; the group is empty if $A=0$. At $K$ take the \emph{differences} of the columns at $K\pm i$, $1\le i\le B$, for even $r$, and at $K\pm(i+1/2)$, $0\le i<B$, for odd $r$. All these columns vanish at the simple nodes, as $\cos(M\theta_l)=\sin(K\theta_l)=0$. These pairs lie in the original interval: in the even case it is $[M-A+1,K+B]$, in the odd case it is $[M-A+1,K+B-1/2]$. They are disjoint since the occupied half-widths fit strictly between $M$ and $K$. In each sum or difference keep one original column. The simple block has full rank: for even $r$ the interval contains $K-t+1,\ldots,K$, and for odd $r$ it contains $K-t+1/2,\ldots,K-1/2$. Their evaluations, up to row factors and column reversal, are respectively those of the first $t$ monic orthogonal polynomials, with the constant column doubled for even $r$. Here use $\sin K\theta_l=0$, $\cos K\theta_l=(-1)^{n+1}$. The absolute value of the complementary simple determinant is constant as $n$ varies.

For the mirror order $rn+r+1-s$ the frequency set is the reflection of the preceding set about $K$, in reverse order. The same operations therefore apply with $M$ replaced by $M'=2K-M$, with the same $A,B$. The centers have gap $+r/2$ or $-r/2$, both fixed and nonzero. Sums at the half-period center transform into sums and differences at $K$ into differences, up to a column sign. Oscillations of the simple rows have already been accounted for.

After expansion on the simple rows and triangular changes of basis, the absolute value is a constant independent of $n$ times the absolute jet determinant of $x^i U$ for $0\le i<A$ and $x^i W$ for $0\le i<B$. Indeed the extra pair factors are, up to constants of fixed absolute value, polynomials of successive degrees: at $M$ they are $1,2\cos(i\phi)$; at $K$ they are $\sin(i\phi)/\sin\phi$ in the even case and $\cos((i+1/2)\phi)/\cos(\phi/2)$ in the odd case. The generators, with $M$ denoting either chosen center, are
\[
\begin{array}{c|cc}
 r&U&W\\ \hline
 2t&\cos(M\phi)&\sin(K\phi)\sin\phi\\
 2t+1&\sin(M\phi)/\sin(\phi/2)&\sin(K\phi)\cot(\phi/2).
\end{array}
\]

\subsection*{Exact degrees}
For odd $r$ both jets are units, with $U(0)=2M$ and $W(0)=2K$. Divide by the generator of the larger group, of size $u=\max(A,B)$, in the ring $\mathbb R[x]/(x^t)$, extracting a factor of order $M^t$. On expansion along the $u$ monomials the residual determinant of size $v=\min(A,B)$ takes coefficients indexed by $h=u+i-j'$, $0\le i,j'<v$. If $v>0$ all $h$ are positive. Up to its value at zero, the ratio of the smaller generator to the larger has at $x=z/H^2$ the expansion
\[
 1+\frac{\delta}{H}(\sqrt z\cot\sqrt z-1)+O(H^{-2}).
\]
Here $H$ is the corresponding larger-group frequency, $K$ or $M$, and $H+\delta$ is the other. To verify this, the sine quotient gives
$\frac{H}{H+\delta}(\cos(\delta\phi)+\cot(H\phi)\sin(\delta\phi))$
and the remaining factor $\cos(\phi/2)^{\pm1}$ equals $1+O(H^{-2})$. This is uniform on a sufficiently small fixed circle in $z$, since
$\phi=2\arcsin(\sqrt z/(2H))$. The coefficients at positive $h$ therefore have leading order $H^{2h-1}$ and, after reversal of columns, their leading determinant is a nonzero shifted Hankel determinant of $c_h=2\sum_{s\ge1}(\pi s)^{-2h}$. Its strict positivity follows by squaring a nonzero polynomial against that infinite discrete measure. All permutation terms have order $H^{2uv-v}$, so the exact degree is $t+2uv-v=(2v+1)u$. This includes $v=0$. For $1\le j\le\lfloor(r+3)/4\rfloor$ use $s=j$ and $s=t+2-j$. These interchange $A,B$ and their mirrors give the other required orders. This proves the first part of Conjecture~22.

For even $r$, $U(0)=1$ and $(W/x)(0)=K$. If $A\ge B$, division by $U$ leaves the $B$-row coefficient determinant of $W/U$. Its expansion at $x=z/M^2$ is $M^{-1}\sqrt z\tan\sqrt z+O(M^{-2})$. Its strictly nonzero leading minor is a shifted Hankel determinant of $b_h=2\sum_{s\ge1}(\pi(s-1/2))^{-2h}$ with shift $A-B+1\ge1$. This gives degree $2AB-B$. If $B>A$, expansion on the constant row eliminates the unshifted $U$ column. Divide all remaining columns by $x$, then factor $W/x$ from the jets of dimension $t-1$, extracting $K^{t-1}$. The $A-1$ remaining rows use $U/(W/x)$, whose scaled expansion is $M^{-1}\sqrt z\cot\sqrt z+O(M^{-2})$. The shift is $B-A+2\ge1$. Its positive-moment determinant gives degree $(t-1)+2B(A-1)-(A-1)=B(2A-1)$. Both cases include the empty residual determinant. This is $(j-1)(2t+1-2j)$ for both prescribed residues. It proves the second part.

\subsection*{Cigler's proposed adjacent sums for the binomial moments}
Cigler's 24 April 2018 comment on MathOverflow question 297870 \cite{MO18} proposed explicit formulas for $S_r(n)=d_r(rn-1)+d_r(rn+2)$, $n>1$; Tyson's answer on 24 July expressly identified these sums as omitted from their paper, arXiv:1807.08330. The exact evaluations are
\[
\begin{aligned}
 S_{2t+1}(n)&=\left(2-\binom{2t+2}{2}\right)(2n+1)^t &&(t\ge1),\\
 S_{2t}(n)&=(-1)^{tn}\left(2-\binom{2t+1}{2}\right) &&(t\ge2),\\
 S_2(n)&=2(-1)^{n+1},\qquad S_1(n)=2 &&(n\ge1).
\end{aligned}
\]
The first two evaluations also hold for $n=1$. The even formula proposed on MathOverflow has the opposite sign, already giving $+8$ in place of $-8$ at $r=4,n=2$.

For the derivation reflect the frequency block at $rn+2$ about $M=(n+1/2)r$. This gives the reverse of the block at $rn-1$. Now $2M\theta_l$ is an odd multiple of $\pi$: reversal of the $2t$ columns and negation of all $t$ simple rows give total sign $(-1)^{\binom{2t}{2}+t}=1$ and leave identical simple blocks after reflection. Use the preceding folding with $s=2$, with $A=t-1$ and $B=1$, for the ranges of the first two formulas. The $A$ kernel columns at $M$ are common to both matrices. With $K'=2M-K=rn$, the other kernel column, an ordered difference of one pair, changes orientation under reflection. After adding the two determinants it is replaced by the difference between the pair differences at $K$ and $K'$. In the even case the difference at $K$ equals $4(-1)^K\sin(K\phi)\sin\phi$. Since $K-K'=r$ is even, subtracting the analogous expression at $K'$ gives
$8(-1)^K U\sin(r\phi/2)\sin\phi$. In the odd case the pair difference at $K$ is $2(-1)^K\sin(K\phi)\cot(\phi/2)$. Now $r$ is odd, so subtracting the pair difference at $K'$ gives
$4(-1)^K U\cos(r\phi/2)\cos(\phi/2)$.
Thus the entire jet block is divisible by $U$ in both cases, and after division its entries are fixed apart from the signs already accounted for. In the even case $U(0)=1$ and the net sign changes by $(-1)^t$ per period. In the odd case $U(0)=2M=r(2n+1)$ and the signs from the simple and jet rows cancel.

For odd $r\ge3$ the continued value at $N=-1$ vanishes because the numerator columns at $-1$ and zero agree. For even $r\ge4$ extend the numerator frequency rule by $P_u=2\cos(u\theta)$ at all integers, keeping the original monic denominator. Its value at $N=-1$ is zero since the columns with frequencies $-1,1$ agree; this extension is only needed for calibration. Since $d_r(2)=\binom{r+4}2-(r+2)^2=2-\binom{r+1}2$, evaluation at $n=0$ gives the claimed constants. Finally, $d_1(N)=1$ and Cigler--Tyson's $d_2(2n)=d_2(2n+1)=(-1)^n$ give the two remaining cases.

\section{Reflection at a zero of a cosine: the even binomial identities}
Use $b_{r,h}=\binom{2h+r}{h}$ for $h\ge0$ and $b_{r,h}=0$ otherwise; write
$d_{r,s}(N)=\det(b_{r,i+j+s})_{0\le i,j<N}$ with $d_{r,s}(0)=1$.
For $k\ge1$, $m\ge-1$, $t=k+m$ and $n\ge0$, the two even identities in Cigler's Conjecture~17 are
\begin{align*}
d_{2t,-m}(2tn)&=(-1)^{tn},\\
d_{2t,-m}(2tn+m+1)&=(-1)^{tn+\binom{m+1}{2}}.
\end{align*}
We prove both. When $t=0$ both sizes are zero. Otherwise put $x=2+w+w^{-1}$ and $L(f)=\operatorname{CT}_w f(x)$. Define $C_s^{*}(x)=w^s+w^{-s}=2\cos(s\theta)$ when $w=e^{i\theta}$. For positive $s$ these are monic polynomials. For $s=0$ the cosine polynomial is $2$ and the monic basis element is $1$. The squared norms of the monic basis for $L$ are $1,2,2,\ldots$ by constant-term extraction.
The monic multiplier $q=x^k C_t^{*}$ realizes the desired moments as $L(x^u q)/2$. Indeed symmetry in $w,w^{-1}$ reduces this expression to the coefficient of $w^{u-m}$ in $(1+w)^{2(u-m)+2t}=(1+w)^{2(u+k)}$. This also gives zero at negative lower indices.

In \eqref{eq:coll} use the first $k$ coefficients at zero and the $t$ roots $\theta_\ell=(2\ell-1)\pi/(2t)$ of $C_t^{*}$. The norms and the $2^{-N}$ from the multiplier give an overall factor $1/2$ at $N\ge1$. Use $C_0^{*}=2$ also in the numerator at $N=0$; the doubled zeroth column there gives exactly the empty determinant after the same factor $1/2$. The denominator always uses the monic basis.

Write $M=(2n+1)t$. For the first progression the frequencies in the numerator are $M-t,\ldots,M+k-1$. Keep $M-t,\ldots,M-1$ and fold the remaining columns into the singleton at $M$ and the sums at $M\pm i$ for $1\le i<k$. For the second progression the interval is $M-k+1,\ldots,M+t$. Keep $M+1,\ldots,M+t$ and carry out the same folds. The condition $k-1\le t$ ensures every partner exists. The folded columns vanish on all simple rows; they are $C_M^{*}$ times monic polynomials of degrees $0,\ldots,k-1$, by $C_{M+i}^{*}+C_{M-i}^{*}=C_M^{*}C_i^{*}$. Their jet determinant depends on $M$ only through $C_M^{*}(0)^k=(2(-1)^M)^k$, which stays fixed when $n$ increases. In the retained $t$ columns the shift by $2t$ negates every simple row. The permutation is fixed and the external sign $(-1)^{N(k+t)}$ also stays fixed. Each determinant consequently changes by $(-1)^t$ per step, including the step from $n=0$. At $n=0$ the initial values are $1$ and $(-1)^{\binom{m+1}{2}}$. For $m\ge0$ the latter is the determinant of a matrix of order $m+1$ with ones on the antidiagonal and zeros above that antidiagonal. For $m=-1$ it is the empty determinant. This proves the formulas.

\section{Odd binomial Hankel determinants in Cigler's Conjecture 17}
Let $k\ge1$, $m\ge-1$, $t=k+m$, $p=2t+1$, and write
$h(N)=d_{p,-m}(N)$ for the binomial determinant defined in
Section~4 of \cite{C23}. Thus $h(0)=1$ and the input
$\binom{2u+p}{u}$ is set to zero for $u<0$.
We prove, for every $n\ge0$,
\begin{align}
 h(pn)&=(2n+1)^k, &
 h(pn+m+1)&=(-1)^{\binom{m+1}{2}}(2n+1)^k.\label{eq:odd17a}
\end{align}
For $m\ge0$ the third proposed identity also holds:
\begin{equation}
 h(pn+t+1)=(-1)^{\binom{t+1}{2}}4^k(n+1)^k.\label{eq:odd17b}
\end{equation}
At $m=-1$ this last identity fails as an identity in $n$ for
\emph{every} $k\ge1$; we give an exact replacement below.
For example, $k=2,m=-1,n=0$ gives
$h(2)=5\cdot84-21^2=-21$ in place of $-16$.

\paragraph{Moment functional.}
Put $x=2+z+z^{-1}$ and $L(f)=[z^0]f(x)$. For $u\in\mathbb Z$ write
$C_u(x)=z^u+z^{-u}$, with $C_0=2$, and put $A_0=1$, $A_u=C_u$ for
$u\ge1$. These $A_u$ are monic and orthogonal for $L$, with squared
norms $1,2,2,\ldots$. Taking the constant term of
$x^u z^m(1+z)^{2k+1}$ gives
$\binom{2(u-m)+p}{u-m}$ for $u\ge0$, using the stated zero convention
when $u<m$. Averaging this Laurent multiplier with its image under
$z\mapsto z^{-1}$ yields $x^{k+1}R_t(x)/2$, where
\[
 R_t(x)=\frac{C_{t+1}(x)+C_t(x)}{x}
       =\frac{\cos((t+1/2)\theta)}{\cos(\theta/2)},
 \qquad x=2+2\cos\theta.
\]
Thus $R_t$ is monic of degree $t$ (with $R_0=1$), and its distinct
nonzero roots correspond to
$\theta=(2\ell+1)\pi/p$, $0\le\ell<t$.

Write $d=t+k+1$ and let $V(f)$ consist of the evaluations at these
$t$ roots followed by the coefficients $[x^i]f$, $0\le i\le k$.
The polynomial-modification determinant identity gives the bilateral
extension
\[
 F(N)=\frac{(-1)^{Nd}}{2}
 \frac{\det(V(C_{N+s}))_{0\le s<d}}
      {\det(V(A_s))_{0\le s<d}}.
\]
The denominator is nonzero by polynomial division. For $N>0$ the
unmodified Gram determinant is $2^{N-1}$ and the scalar
in the multiplier contributes $2^{-N}$. Hence $F(N)=h(N)$.
At $N=0$ precisely its first column is doubled relative to the
denominator, which gives $F(0)=1$ as well. The evenness $C_{-u}=C_u$
reverses the columns, giving
$F(1-d-N)=(-1)^{\binom d2}F(N)$.

\paragraph{A half-integer center.}
For $N=pn$ keep the first $t+1$ columns, with indices
$pn,\ldots,pn+t$. For each $1\le i\le k$ add the column with index
$pn+t+1-i$ to the column with index $pn+t+i$. The companion belongs
to the retained block since $k\le t+1$. The resulting $k$ columns
vanish at the simple roots because their frequencies are symmetric
about $M=p(n+1/2)$. At $x=2-2\cos\phi$ they equal
\[
 -4(-1)^{pn+t+i}\sin(M\phi)\sin((i-1/2)\phi).
\]
They also vanish at $x=0$. Divide these columns by
$V_1=\sin(M\phi)\sin(\phi/2)$. The quotient factors
$\sin((i-1/2)\phi)/\sin(\phi/2)=1+2\sum_{s=1}^{i-1}\cos(s\phi)$
are polynomials of degree $i-1$ in $x$ with leading coefficient
$(-1)^{i-1}$. As $[x]V_1=M/2$, the block of coefficients of
degrees $1,\ldots,k$ has determinant
$[2M(-1)^{pn+t}]^k$.
Let $U$ denote the remaining block from $C_0,\ldots,C_t$ at the
simple rows and the constant row. For general $n$ this block is
$(-1)^n U$, because both the simple-root angles and $\theta=\pi$
acquire that sign under a shift of frequency by $pn$.
At $n=0$, $F(0)=1$ ensures that $U$ is nonsingular. The signs from
the two blocks and $(-1)^{Nd}$ cancel, so on all integers $n$,
$F(pn)=(2n+1)^k$. Since $1-d-(pn+m+1)=-p(n+1)$, the reflection
gives the second formula in \eqref{eq:odd17a} using
$\binom d2+k\equiv\binom{m+1}2\pmod2$.

\paragraph{An integer center.}
For $N=pn+t+1$ set $K=p(n+1)$. The column indices are
$K-t,\ldots,K+k$. Keep $K-t,\ldots,K$; their simple and constant
rows give $(-1)^{n+1}U$ with the column order reversed. For
$m\ge0$, $k\le t$, so one may subtract the retained column $K-i$
from $K+i$ for each $1\le i\le k$. These differences vanish at
the simple roots and have the form
$-4(-1)^{K+i}\sin(K\phi)\sin(i\phi)$ at the zero jets.
Here $V_2=\sin(K\phi)\sin\phi$ has $[x]V_2=K$, and
$\sin(i\phi)/\sin\phi$ is of degree $i-1$ in $x$ with leading
coefficient $(-1)^{i-1}$. Thus the new jet determinant equals
$[4K(-1)^K]^k$. Relative to the calibration $F(0)=1$, the
absolute ratio is $(4K/p)^k$. Its sign exponent modulo two is
\[
 Nd+(n+1)(t+1)+\binom{t+1}{2}+(K-t)k
 \equiv\binom{t+1}{2},
\]
which proves \eqref{eq:odd17b}.

\paragraph{Exact replacement at $m=-1$.}
Here $k=t+1$. Use the same first $t$ differences, then add the
column $K-t$ to $K+t+1$. Their indices differ by $p$, so this last
sum still vanishes on the simple rows; at the zero jets it equals
\[
 -4(-1)^{K+t+1}\sin((K+1/2)\phi)\sin((t+1/2)\phi).
\]
The first $t$ difference columns span $V_2$ times polynomials
of degrees $0,\ldots,t-1$. After removal of the common factor
$V_2$ in the shifted jets (equivalently, division by the unit
$V_2/x$ after removing $x$), the one remaining entry is the
coefficient of $x^t$. Consequently the exact answer is the
right side of \eqref{eq:odd17b} multiplied by
\[
 (-1)^t [x^t]\,
 \frac{\sin((K+1/2)\phi)\sin((t+1/2)\phi)}{\sin(K\phi)\sin\phi},
 \qquad x=2-2\cos\phi .
\]
All quotients admit Taylor expansions at the origin for $K>0$.
At $t=0$ the corrected value is $2n+3$. For $t\ge1$ the
coefficient in this display, before its external $(-1)^t$, is
asymptotic to
\[
 \frac{t+1/2}{2}\,K^{2t-1}[y^t](\sqrt y\cot\sqrt y).
\]
Indeed the first sine ratio equals
$\cos(\phi/2)+\cot(K\phi)\sin(\phi/2)$, and the second ratio
is $(t+1/2)+O(\phi^2)$. Substituting $x=y/K^2$, the constant
term in $y$ at order $K^0$ is $t+1/2$, and the term of order
$K^{-1}$ is $(t+1/2)\sqrt y\cot\sqrt y/2$.
These expansions hold uniformly on a small fixed circle in $y$.
The coefficient displayed is nonzero by the partial fractions of
the cotangent. The actual asymptotic growth has exponent
$k+2t-1=3k-3>k$ for $k\ge2$. This proves the claimed failure
for every parameter on this boundary.

\section{One-variable identities for the cubic convolution}\label{sec:cubicprime}
For a prime $p>3$ and $x\in\mathbb Z_p$, put
\[
 T_p(x)=\sum_{k=0}^{p-1}(-1)^k(2k+1)\sum_{j=0}^k
          \binom{-x}{j}^{\!3}\binom{x-1}{k-j}^{\!3}.
\]
The third assertion of Sun's Conjecture 35(v) \cite[p.~31]{Sun} concerns exactly $T_p(1/4)$. We first derive the identities needed at arbitrary parameters.


For complex $z$ in a neighborhood of a nonnegative integer, set
\[
 F(z)=\sum_{j\geq0}(-1)^j\binom zj^3,\qquad
 K(z)=\sum_{j\geq0}(2j-z)(-1)^j\binom zj^3.
\]
The Dixon and Dougall summations give
\begin{align}
 F(z)&=\cos(\pi z/2)\frac{\Gamma(1+3z/2)}{\Gamma(1+z/2)^3},\label{eq:suncubic-dixon}\\
 K(z)&=-2\sin(\pi z/2)\frac{\Gamma((1+3z)/2)}{\Gamma((1+z)/2)^3}.\label{eq:suncubic-dougall}
\end{align}
Indeed, take $a=b=c=-z$ in Dixon's $\mathstrut{}_3F_2(1)$ summation.
In Dougall's $\mathstrut{}_5F_4(1)$ summation, with upper parameters
$a,1+a/2,b,c,d$, take $a=b=c=-z$ and $d=(1-z)/2$.
The last parameter cancels the lower parameter $1+a-d$. Multiplying the sum by $-z$ introduces the factor $2j-z$.
Euler's reflection formula gives the displayed expressions.
Both summation formulas are recorded together by Masjed-Jamei and Koepf
\cite[eqs.~(8), (13)]{masjedkoepf}.
The parameter substitutions hold off the isolated singularities of the hypergeometric notation and extend there by continuity.
For $\Re z>-1/3$ both series converge locally uniformly. Indeed,
$(-1)^j\binom zj=\Gamma(j-z)/(\Gamma(-z)\Gamma(j+1))$ is
$O(j^{-1-\Re z})$ locally uniformly for large $j$, by the gamma-ratio estimate.

For $N\ge0$ and $j>N$, $\binom zj^3$ has a zero of order three at $z=N$.
Consequently the values and first two derivatives of $F,K$ at $N$ agree with those of their sums through $j=N$.
In particular these derivatives are rational numbers, obtained by differentiating finite polynomials.
The gamma formulas yield $K(2u)=F(2u+1)=0$ and
\begin{align*}
 F(2u)&=(-1)^u\frac{(3u)!}{u!^3},&
 K'(2u)&=-(-1)^u\frac{(6u)!\,u!^3}{(3u)!\,(2u)!^3},\\
 K(2u+1)&=-2(-1)^u\frac{(3u+1)!}{u!^3},&
 F'(2u+1)&=-2(-1)^u\frac{(6u+4)!\,(u+1)!^3}{(3u+2)!\,(2u+2)!^3}.
\end{align*}
These follow directly from $\Gamma(t+1/2)=(2t)!\sqrt\pi/(4^t t!)$ for integers $t\geq0$.

\subsection*{Arithmetic of complementary rows}
Put $X=-x$, $Y=x-1$, and take $A,B\in[0,p-1]$ representing $X,Y$
modulo $p$. Thus $A+B=p-1$ and $A,B$ have the same parity. Write
$X=A+p\sigma$, $Y=B+p\tau$ with $\sigma,\tau\in\mathbb Z_p$.
All binomials of degree less than $p$ have $p$-integral coefficients.
For $A<j<p$, $\binom Xj^3\equiv0\pmod{p^3}$, and the same applies to the second row beyond $B$.
The surviving rectangle $0\leq i\leq A$, $0\leq j\leq B$ lies in $i+j<p$.
Since $2(i+j)+1=(2i-X)+(2j-Y)$, its weighted double sum factors into two products,
the centered moment in either row times the zeroth moment in the other.
For the subsequent Taylor expansion denote the linear coefficient contributed by perturbing the first row, before multiplication by $p\sigma$, by
\[
 C(A,B)=\begin{cases}K'(A)F(B),& A\text{ even},\\F'(A)K(B),& A\text{ odd}.\end{cases}
\]
The role of $A$ and $B$ interchanges in the other product.

\section{A prime-cube evaluation and a counterexample in part 35(v)}
\label{sec:cubic-p3}
Keep the notation $X=A+p\sigma=-x$ and $Y=B+p\tau=x-1$ from Section~\ref{sec:cubicprime}, so $\sigma+\tau=-1$. For every prime $p>3$ and $3x\not\equiv1,2\pmod p$ the following stronger formula holds:
\begin{equation}\label{eq:cubicsunprime3}
 \frac{T_p(x)}{p^2}\equiv R(A,\sigma)+R(B,\tau)\pmod p,\qquad
 R(A,\sigma)=
 \begin{cases}
 -\sigma(1+3\sigma)/(3A+1)& A\ {\rm even},\\
 \phantom{-}\sigma(1+3\sigma)/(3A+2)& A\ {\rm odd}.
 \end{cases}
\end{equation}
Both denominators that occur in \eqref{eq:cubicsunprime3} are units.
We give the second-order calculation, since retaining just the first derivative of the vanishing row misses a term of the same $p$-adic order.

Write $h=(p-1)/2$ and $D(j)=\prod_{i=1}^{j}(2i-1)$, with $D(0)=1$.
The factor of the centered product which vanishes at $A$ is $K$ when $A$ is even and $F$ when $A$ is odd. Let $V(A,B)$ denote the product consisting of that factor in row $X$ and the nonvanishing factor in row $Y$, using the finite row polynomials of degrees indexed through $A$ and $B$. Then $T_p(x)\equiv V(A,B)+V(B,A)\pmod{p^3}$. The derivative coefficient $C=C(A,B)$ from Section~\ref{sec:cubicprime} takes the common form
\[
 C=g\,\frac{D(r)(3h-r)!}{\{D(t)(h-t)!\}^3},
 \qquad
 \begin{array}{c|ccc}
 A&t&r&g\\ \hline
 2u&u&3t&-(-1)^h\\
 2u+1&u+1&3t-1&-2(-1)^h .
 \end{array}
\]
In the generic range $r\ne h$. The numerator has exactly one factor $p$. It comes from $(3h-r)!$ when $r<h$, and from $D(r)$ when $r>h$.

Set $\eta=1$ for even $A$ and $\eta=-1$ for odd $A$. We claim
\begin{equation}\label{eq:cubicsun-unit}
 \frac C p \equiv \frac{\eta}{r-h}
 \begin{cases}1&r<h,\\ -1/2&r>h\end{cases}\pmod p .
\end{equation}
For $0\le j\le h$, reduction of the odd factors gives
$D(j)\equiv(-2)^j h!/(h-j)!\pmod p$.
If $r<h$, then $3h-r=p+(h-1-r)$ and hence
\[
 D(r)(3h-r)!/p\equiv-(-2)^r h!/(h-r)\pmod p.
\]
If $r=h+1+s>h$, then
$D(r)/p\equiv(-2)^h h!2^s s!\pmod p$ and
$(3h-r)!=(p-2-s)!\equiv(-1)^s/(s+1)!\pmod p$.
Their product is $(-2)^{r-1}h!/(r-h)$ modulo $p$.
Dividing in both cases by $((-2)^t h!)^3$, inserting $g$ and using $(h!)^2\equiv(-1)^{h+1}\pmod p$ proves \eqref{eq:cubicsun-unit}.

We now compute the quadratic Taylor terms. If $J$ is the analytic function vanishing at $A$, the logarithmic derivative of $J(z)/(z-A)$ at $A$ is
\[
 L_A=3\sum_{j=t+1}^{r}\frac{1}{2j-1}.
\]
The logarithmic derivative at $B$ of the other factor is
\[
 L_B=\frac32\sum_{j=h-t+1}^{3h-r}\frac1j.
\]
These formulas follow from \eqref{eq:suncubic-dixon}--\eqref{eq:suncubic-dougall} and $\psi(z+1)-\psi(z)=1/z$ for the logarithmic derivative $\psi=\Gamma'/\Gamma$. The derivative of the residual trigonometric factor vanishes in each case. The interpretations for empty sums are zero.
As each omitted binomial cube has a triple zero, the analytic and finite Taylor coefficients agree through order two. The finite polynomials have $p$-integral coefficients. Their expansion therefore gives
\[
 V(A,B)\equiv p\sigma C\bigl(1+p\sigma L_A+p\tau L_B\bigr)\pmod{p^3}.
\]
When $r<h$, the only nonintegral harmonic summand is at $j=p$ in $L_B$. Thus $p L_A\equiv0$ and $pL_B\equiv3/2\pmod p$. When $r>h$, only $2j-1=p$ in $L_A$ contributes, and $pL_A\equiv3$, $pL_B\equiv0\pmod p$. Since $1+3\tau/2=-(1+3\sigma)/2$, both cases of \eqref{eq:cubicsun-unit} yield
\[
 V(A,B)/p^2\equiv-\frac{\eta\,\sigma(1+3\sigma)}{2(r-h)}\pmod p.
\]
For even $A$ the denominator is $3A-p+1$; for odd $A$ it is $3A-p+2$. Adding the reversed contribution proves \eqref{eq:cubicsunprime3}.

\subsection*{The quarter parameter}
For $p>5$, substituting $x=1/4$ in \eqref{eq:cubicsunprime3} gives
\[
 T_p(1/4)\equiv
 \begin{cases}
 p^2 & p\equiv1,7\pmod 8,\\
 -19p^2/5& p\equiv3,5\pmod 8
 \end{cases}\qquad(\bmod\ p^3).
\]
For verification, the respective choices in the two summands are
\[
\begin{array}{c|c|cc|c}
 p\bmod8&A,B\ \text{parity}&\sigma&\tau&
 R(A,\sigma)+R(B,\tau)\\ \hline
 1&\text{even}&-1/4&-3/4&1\\
 3&\text{even}&-3/4&-1/4&-19/5\\
 5&\text{odd}&-1/4&-3/4&-19/5\\
 7&\text{odd}&-3/4&-1/4&1
\end{array}
\]
Sun's printed third assertion in Conjecture 35(v) specifies the value $p^2$ for $p\not\equiv5\pmod8$. It fails for every prime $p>3$ with $p\equiv3\pmod8$, since $-19/5-1=-24/5$ is a unit for such $p$. At $p=11$ one can verify the failure without any summation theorem. The eleven terms, indexed by the outer summation from $k=0$ through $10$, reduce modulo $1331$ to
\[
 1,\ 750,\ 643,\ 1275,\ 667,\ 759,\ 499,\ 496,\ 1073,\ 119,\ 978.
\]
Their sum is $605$ modulo $1331$, whereas $11^2=121$.
For the small primes, direct substitution gives
$T_3(1/4)=61893/16384\equiv9\pmod{27}$ and
$T_5(1/4)=30165631255/4294967296\equiv30\pmod{125}$.
The latter prime is excluded in the original clause. Among its admissible primes the proposed value therefore holds exactly for $p=3$ and for $p\equiv1,7\pmod8$.

\section*{Checks and historical scope}
The polynomial degree and coefficient formulas were checked independently by direct elimination in the original binomial-coefficient matrices modulo $1{,}000{,}003$. In 16 even cases and 18 odd residue cases, the difference of order $d$ equalled $d!$ times the predicted coefficient and the next difference vanished, where $d$ is the asserted degree. The largest even matrix had order 784 and the largest odd matrix had order 342. For example the signed integer values for $D_{7,0}(7n+2)$ at $n=0,1,2,3$ are $-14,315,1673,4746$. Their third difference is $686$, giving the cubic coefficient $343/3$. A separate exact-arithmetic check covered the three odd binomial identities, including the boundary replacement, at $1\le k\le4$, $-1\le m\le2$, $0\le n\le2$. For Wang--Xin 22, a second calculation used the original binomial matrices at $r=1,\ldots,18$ and $27,\ldots,34$, with order at most $1000$. At 205 tested progressions the top difference was nonzero and all 6276 tested differences of the next order vanished modulo $1{,}000{,}003$. The same calculation checked 1037 adjacent sums against their signed formulas. These finite tests are checks on indexing and signs. The proofs supply the unbounded quantifiers.

The original numbered targets and their definitions are given in \cite{C18,C23,Sun}. For the 2018 constant sum the original has the same odd power in both terms. The restatement in Wang--Xin's Conjecture 7 and the corresponding MathDB problem display use different powers. The posted MathDB calculation identifies that misprint. That page's editorial status attributes settlement to Fulmek, while its posted calculation expressly leaves the original identity unproved \cite{MathDB}. Fulmek's stated main theorem and Cigler's subsequent proof compare different backward shifts at fixed order \cite{Back}. Chern--Shi also use Chebyshev reductions and cover several special residues; Han and Han--Yang treat further families of shifted and Narayana convolution determinants \cite{CS,Han,HY}. The public formal proof of Cigler's 2023 Conjecture 11 already supplies the positive central-power identity in 2018 Conjecture 7.2 upon choosing its shift parameter $k-1$ \cite{TT}. The $k=1$ exception of 2018 Conjecture 7.5 was previously recorded \cite{TT}. Cigler--Tyson give the $m=0$ binomial evaluations; Chern--Shi include further overlap when $k=1$, including the boundary value $d_{1,1}(N)=2N+1$ already implicit in \cite[Cor.~18]{KY}. These cases are included for a complete classification.

The original papers, the subsequent papers and problem records cited here, and searches for the reflected constant sums, the two polynomial degree formulas of Wang--Xin 22, the signed shift law, and the quarter-parameter congruence did not identify an earlier proof covering the full results highlighted above. This is a bounded literature assessment as of 20 November 2026. Historical priority of the general arguments and independent external validation remain open. No ranking of these theorems against the entire corpus has been established.

\begin{thebibliography}{99}
\bibitem{C18} J. Cigler, \emph{Catalan numbers, Hankel determinants and Fibonacci polynomials}, 2018, v3, Conjectures 7.2 and 7.5. \href{https://arxiv.org/abs/1801.05608v3}{arXiv:1801.05608v3}.
\bibitem{WX} Y. Wang and G. Xin, Hankel determinants for convolution powers of Catalan numbers, \emph{Discrete Math.} \textbf{342} (2019), 2694--2716. \href{https://doi.org/10.1016/j.disc.2019.05.025}{doi:10.1016/j.disc.2019.05.025}.
\bibitem{CT} J. Cigler and M. Tyson, An interesting class of Hankel determinants, in \emph{Algorithmic Combinatorics}, Springer, 2020, 137--157. \href{https://arxiv.org/abs/1807.08330}{arXiv:1807.08330}.
\bibitem{C23} J. Cigler, \emph{Some experimental observations about Hankel determinants of convolution powers of Catalan numbers}, 2023, v2. \href{https://arxiv.org/abs/2308.07642v2}{arXiv:2308.07642v2}.
\bibitem{MO18} J. Cigler and M. Tyson, comment of 24 April and answer of 24 July 2018, \emph{Some more binomial coefficient determinants}, \href{https://mathoverflow.net/questions/297870}{MathOverflow question 297870}.
\bibitem{MO} J. Cigler, Hankel determinants for some convolutions of Catalan numbers, MathOverflow, March 2023. \href{https://mathoverflow.net/questions/442285}{Question 442285}.
\bibitem{Krat} C. Krattenthaler, Hankel determinants of linear combinations of moments of orthogonal polynomials, II, \emph{Ramanujan J.} \textbf{61} (2023), 597--627. \href{https://arxiv.org/abs/2101.04225}{arXiv:2101.04225}.
\bibitem{KY} C. Krattenthaler and D. Yaqubi, Some determinants of path generating functions, II, 2018. \href{https://arxiv.org/abs/1802.05990}{arXiv:1802.05990}.
\bibitem{Back} M. Fulmek, Hankel determinants of convoluted Catalan numbers and nonintersecting lattice paths, \href{https://arxiv.org/abs/2402.19127}{arXiv:2402.19127}. J. Cigler, Hankel determinants of convolution powers of Catalan numbers revisited, \href{https://arxiv.org/abs/2403.11244}{arXiv:2403.11244}, 2024.
\bibitem{CS} S. Chern and W. Shi, \emph{Hankel determinants of Catalan-like sequences}, 2026. \href{https://arxiv.org/abs/2608.27208}{arXiv:2608.27208}.
\bibitem{Han} G.-N. Han, Hankel determinants for convolution powers of Narayana polynomials, \emph{Electron. J. Combin.} \textbf{33} (2026), P2.42. \href{https://doi.org/10.37236/14995}{doi:10.37236/14995}.
\bibitem{HY} G.-N. Han and L. Yang, Shifted Hankel determinants via Hankel continued fractions, 16 November 2026. \href{https://arxiv.org/abs/2611.21968}{arXiv:2611.21968}.
\bibitem{TT} The Omega Institute, \href{https://www.trureturing.org/research/cigler-catalan-power-shifted-hankel/}{formal proof dossier for Cigler Conjecture 11}, and \href{https://www.trureturing.org/knowledge/node/c875051dc1dd7194f09bc1f09b598c9e682cdba01056d248b04dfcb03c3a82c1/}{dossier on the printed endpoint of Conjecture 7.5}.
\bibitem{MathDB} MathDB, \href{https://mathdb.com/p/340916/cigler-s-odd-index-hankel-determinant-formulas-for-catalan-c}{problem 340916}, editorial status and posted source correction, accessed 20 November 2026.
\bibitem{Sun} Z.-W. Sun, Open conjectures on congruences, \emph{Nanjing Univ. J. Math. Biquarterly} \textbf{36} (2019), 1--99, Conjecture 35(v), p.~31. \href{https://doi.org/10.3969/j.issn.0469-5097.2019.01.01}{doi:10.3969/j.issn.0469-5097.2019.01.01}.
\bibitem{masjedkoepf} M. Masjed-Jamei and W. Koepf, Some summation theorems for generalized hypergeometric functions, \emph{Axioms} \textbf{7} (2018), 38. \href{https://doi.org/10.3390/axioms7020038}{doi:10.3390/axioms7020038}.
\bibitem{VMRC} QEEDDD, \emph{Verifiable Mathematical Reasoning Corpus}, \href{https://github.com/QEEDDD/Verifiable-Mathematical-Reasoning-Corpus/tree/c2ac34ea4622778187173f41bb1996deee8081b7}{snapshot c2ac34e}, accessed 20 November 2026, especially manuscripts 026 and 079.
\end{thebibliography}
\end{document}
